\cotefile{46}
\scanpage{65}
% archive sheet number: 64
% L1-5
$a\in\unsure{\mathscr{B}|\mathcal{B}}(X)$, soit $a^\circ$ la classe ``opposée'' (on remplace une algèbre par l'algèbre opposée), alors $aa^\circ$ est trivial. En effet, on a canoniquement $A\otimes A^\circ\simeq L(A,A)$ : D'où la définition du \uline{groupe de Brauer}, noté $\operatorname{Br}(X)$.

% L6-9
On \unsure{note|voit} qu'on a
\[
\delta(ab)=\unsure{I_m|\iota_m}\delta(a)+\unsure{I_n|\iota_n}\delta(b)
\]
\margin{compte tenu des hom
$\begin{gathered}\mathbb{Z}/n\mathbb{Z}\\\mathbb{Z}/m\mathbb{Z}\end{gathered}\longrightarrow\mathbb{Z}/mn\mathbb{Z}$}
d'où résulte qu'on a un hom de \deleted{\unsure{groupes|groupe}} \inserted{monoïdes}
\[
\mathscr{B}(X)\longrightarrow H^2(X,\mathbb{Q}/\mathbb{Z})
\]
% L10-14
Cet homomorphisme définit, par composition
$\unsure{B|\mathscr{B}}(X)\longrightarrow H^2(X,\underline{\mathbb{C}}^*)$, \deleted{\unsure{d'où}} et \deleted{\unsure{est trivial}} \inserted{son noyau est} \deleted{\unsure{sur les}\ \illegible{2}} \inserted{formé des} éléments ``triviaux''. \deleted{\unsure{on}} D'où un hom naturel \uuline{injectif}
\[
\boxed{\operatorname{Br}(X)\xrightarrow{\unsure{\varphi|\psi}}H^2(X,\underline{\mathbb{C}}^*)}
\]
% L15-16
dont l'image est contenue dans $\operatorname{Im}H^2(X,\mathbb{Q}/\mathbb{Z})$, donc dans le sous-groupe de Torsion.

% L17-20
\uline{Question 1} Tout élément de torsion de $H^2(X,\underline{\mathbb{C}}^*)$ provient-il d'un élément de $\operatorname{Br}(X)$ ?

\unsure{Cela|Ce} résulterait de \unsure{l'affirmative|l'affirmation} à la question suivante plus précise

% L21-22
\uline{Question 2} Tout élément de $H^2(X,\mathbb{Z}/n\mathbb{Z})$ est-il de la forme $\delta a$, avec $a\in H^1(X,\unsure{\operatorname{GP}|\operatorname{Gp}}(n,\ ))$ ???

% The source leaves the second GP argument blank.
